\documentclass[12pt]{article}
\usepackage{amsmath}
\usepackage{amssymb}
\usepackage{geometry}
\geometry{a4paper, margin=1in}

\title{\textbf{Explicit Construction of the Asymptotic Trapping Region}}
\date{}

\begin{document}

\maketitle

\section{Introduction}
To complete Part 3 of the computer-assisted proof, we must mathematically construct a forward-invariant trapping region $\mathcal{T}$ that captures the numerical flow at large radius $r = R$ and rigorously forces the state into the exact asymptotic cone as $r \to \infty$.

\section{The Ratio Variables and Vector Field}
We define bounded ratio variables designed to limit to constants based on the proven strictly conical asymptotics:
\[ X = \frac{r a'}{a}, \quad Y = \frac{r b'}{b}, \quad W = \frac{f'}{r}, \quad B = \frac{b}{r} \]
The target limits as $r \to \infty$ are $(X, Y, W) \to (1, 1, -1)$ while $B \to b_\infty > 0$.

By differentiating with respect to $r$ and substituting the original ODEs ($a'', b'', f''$), we derive the autonomous-like vector field (using $\tau = \ln r$ such that $\frac{d}{d\tau} = r \frac{d}{dr}$):
\begin{align*}
\frac{dX}{d\tau} &= X - X^2 - 2 \frac{X Y}{r} + X W r^2 + r^2 \\
\frac{dY}{d\tau} &= Y - Y^2 + \frac{1}{r^2 B^2} - \frac{Y^2}{r^2} - \frac{X Y}{r^2} + Y W r^2 + r^2 \\
\frac{dW}{d\tau} &= -W + \frac{1}{r^2} \left[ -2 \frac{X Y}{r} + X W r^2 + r^2 + 2\left(\frac{1}{r^2 B^2} - \frac{Y^2}{r^2} - \frac{X Y}{r^2} + Y W r^2 + r^2\right) - r^2 \right] \\
\frac{dB}{d\tau} &= B (Y - 1)
\end{align*}
Note the critical cancellation: as $X \to 1$ and $W \to -1$, the dominant diverging terms $X W r^2 + r^2 \approx -r^2 + r^2 = 0$ perfectly cancel, stabilizing the vector field.

\section{The Boundary Inequalities}
We define the trapping region $\mathcal{T}$ at $r=R$ as a hyper-box around the limits:
\[ \mathcal{T}(r) = \{ |X - 1| \le \delta_X, \quad |Y - 1| \le \delta_Y, \quad |W + 1| \le \delta_W \} \]
For $\mathcal{T}$ to be forward-invariant, the vector field on the boundary faces must point strictly inward. This translates to the following mathematical inequalities evaluated at $r \ge R$:

\begin{enumerate}
    \item \textbf{On the $X$-boundaries} ($X = 1 \pm \delta_X$):
    \[ \pm \left( X - X^2 - 2 \frac{X Y}{r} + X W r^2 + r^2 \right) < 0 \]
    
    \item \textbf{On the $Y$-boundaries} ($Y = 1 \pm \delta_Y$):
    \[ \pm \left( Y - Y^2 + \frac{1}{r^2 B^2} - \frac{Y^2}{r^2} - \frac{X Y}{r^2} + Y W r^2 + r^2 \right) < 0 \]
    
    \item \textbf{On the $W$-boundaries} ($W = -1 \pm \delta_W$):
    \[ \pm \left( \frac{dW}{d\tau} \right) < 0 \]
\end{enumerate}

\section{Computational Verification}
To complete the proof in code, one chooses explicit constants for the widths (e.g., $\delta_X = 0.05, \delta_Y = 0.05, \delta_W = 0.01/R$). 
The interval bounding box generated by the Julia script at $r=R$ is projected into the variables $(X, Y, W, B)$. If this projected box lies completely inside $\mathcal{T}(R)$, and interval arithmetic confirms the above boundary inequalities evaluate to strictly less than $0$ over the semi-infinite domain $[R, \infty)$, the rigorous proof of asymptotic convergence is finalized.

\end{document}
